\chapter{Back Injuries: Symptoms and Treatments}
\index{Back Injuries: Symptoms and Treatments}

\section*{Definition and Overview}
Back injuries are injuries to the bones, joints, muscles, ligaments, discs, or nerves of the back. They range from a simple muscle strain that settles in a few days to serious injuries such as a fractured vertebra or spinal cord damage. The lower back is the area most often injured, followed by the neck and the mid-back. Most back injuries are minor and heal with self-care, but a small number need urgent assessment because they involve the spinal cord or nerve roots.

An understanding of what has been injured is important, because the treatment for a muscle strain is very different from that for a disc prolapse or a fractured spine. Red-flag features, such as loss of bladder or bowel control, weakness, or numbness in the legs, always require urgent medical assessment. A careful history and examination, together with selective imaging, allow the clinician to distinguish the common benign mechanical causes of back pain from the small number of serious causes.

\section*{Epidemiology}
Back pain and injury are extremely common. Around 1 in 6 people has low back pain at any given time, and most adults experience at least one episode of back pain in their lifetime. Acute back injury is common in people doing heavy manual work, in sport, and after falls or car crashes. Disc problems and arthritic changes become more common with age, and osteoporotic (fragility) fractures of the spine mainly affect older women. Back pain is one of the most frequent reasons for visiting a GP and for time off work, and it is among the leading causes of disability worldwide.

\section*{Aetiology and Pathophysiology}
\begin{itemize}
    \item \textbf{Muscle and ligament strain:} heavy lifting, twisting, sudden or awkward movements, and poor lifting technique -- the most common cause
    \item \textbf{Disc injury:} a disc bulge or herniation pressing on a nerve root, causing radiating pain (for example, sciatica)
    \item \textbf{Vertebral fracture:} high-energy trauma in younger people (falls, car crashes) or low-energy fractures through osteoporotic bone in older adults
    \item \textbf{Joint and degenerative change:} wear and tear (spondylosis), facet joint irritation, and spinal stenosis
    \item \textbf{Spondylolisthesis:} slippage of one vertebra forward over the one below it
    \item \textbf{Spinal cord and nerve injury:} compression or damage causing weakness, numbness, or loss of function
    \item \textbf{Less common causes:} infection of the spine (discitis, osteomyelitis), inflammatory arthritis, or tumours, which can all present with back pain
\end{itemize}

At a physiological level, the intervertebral discs act as shock absorbers between the vertebrae. With age or repeated strain the outer ring of the disc can tear, allowing the soft inner core to bulge or herniate against nearby nerve roots. Narrowing of the spinal canal (stenosis) compresses the spinal cord or the bundle of nerves known as the cauda equina, which produces pain, numbness, and weakness in a distribution that depends on the level affected.

\section*{Clinical Features}
\begin{itemize}
    \item Pain localised to the back, often worse with movement or certain positions
    \item Stiffness and muscle spasm, with reduced range of motion
    \item Radiating pain, for example sciatica (pain down the back of the leg), tingling, or numbness
    \item Weakness in a leg or foot, or difficulty lifting the foot (foot drop)
    \item Loss of feeling in the saddle area (groin and inner thighs) or loss of bladder or bowel control -- red flags
    \item Tenderness over the spine after a fall or other trauma
    \item In older adults: sudden back pain after minor strain or a fall that may indicate a fragility fracture
\end{itemize}

\section*{Differential Diagnosis}
\begin{itemize}
    \item Muscle strain or sprain versus disc herniation versus fracture
    \item Sciatica from a disc versus compression of nerves elsewhere (for example, hip or piriformis problems)
    \item Spinal stenosis versus vascular claudication (leg pain from poor circulation)
    \item Inflammatory arthritis (for example, ankylosing spondylitis) -- pain and stiffness worse at rest
    \item Infection of the spine (discitis, osteomyelitis) -- fever and severe constant pain
    \item Malignancy (metastatic bone disease, myeloma) -- night pain, weight loss, history of cancer
    \item Non-spinal causes of back pain: kidney stones, pancreatitis, abdominal aortic aneurysm, and pelvic problems
    \item Cauda equina syndrome -- a surgical emergency with loss of bladder or bowel control
\end{itemize}

\section*{When to Seek Medical Care}
Most acute back injuries improve within a few weeks with self-care, but some need urgent attention.

\textbf{Contact your local emergency services for:}
\begin{itemize}
    \item Loss of control of bladder or bowels (or difficulty passing urine)
    \item Numbness in the saddle area or sudden weakness of both legs
    \item Inability to walk or stand after an injury
    \item Suspected spinal injury after a major fall, car crash, or diving accident
\end{itemize}

See a GP promptly if:
\begin{itemize}
    \item Back pain follows a fall or other trauma
    \item Pain is severe, progressive, or present at night
    \item There is numbness, tingling, or weakness in a leg
    \item You have a fever, unintentional weight loss, or a history of cancer
    \item Pain does not improve after a few weeks of self-care
\end{itemize}

\section*{Diagnosis and Investigations}
\begin{itemize}
    \item \textbf{History:} mechanism of injury, location and type of pain, radiation, bladder and bowel function, previous episodes, and osteoporosis risk
    \item \textbf{Examination:} inspection and palpation of the spine, range of movement, posture, and muscle spasm
    \item \textbf{Neurological assessment:} strength, sensation, and reflexes in the legs; straight leg raise test for nerve root irritation
    \item \textbf{Imaging:} X-ray for suspected fracture; CT for detailed bone views; MRI for discs, nerve roots, spinal cord, and soft tissue
    \item \textbf{Blood tests:} for suspected infection or inflammation (full blood count, inflammatory markers, blood cultures)
    \item \textbf{Urgent imaging:} when red-flag features or cauda equina syndrome are suspected
\end{itemize}

\section*{Management}
\begin{itemize}
    \item \textbf{Keep active:} graded activity and normal movement as tolerated; prolonged bed rest is not helpful
    \item \textbf{Pain relief:} simple analgesics such as paracetamol and anti-inflammatory medicines (NSAIDs), heat packs, and gentle stretching
    \item \textbf{Physiotherapy:} exercises to strengthen the core and back muscles, mobilisation, and advice on returning to activity
    \item \textbf{Manual therapy:} spinal manipulation or massage by appropriately qualified practitioners
    \item \textbf{Specific treatment:} short courses of medicines for nerve pain or muscle spasm where indicated
    \item \textbf{Surgery:} considered for cauda equina syndrome, severe or progressive nerve compression, unstable fractures, or structural problems that fail conservative care
    \item \textbf{Rehabilitation:} graded return to work, sport, and daily activities to reduce the chance of recurrence
\end{itemize}

\section*{Complications}
\begin{itemize}
    \item Chronic or recurrent back pain
    \item Persistent nerve compression causing lasting weakness, numbness, or wasting of leg muscles
    \item Cauda equina syndrome with permanent loss of bladder or bowel control if not treated urgently
    \item Osteoporotic vertebral fractures leading to height loss, deformity, and reduced mobility
    \item Reduced work capacity, physical deconditioning, and depression
    \item Complications of surgery: infection, bleeding, nerve damage, or failure to relieve pain
\end{itemize}

\section*{Prognosis and Outcomes}
Most episodes of acute low back pain settle within four to six weeks, regardless of whether a clear injury occurred. Recurrence is common, so strengthening exercises and good lifting habits reduce the chance of further episodes. Sciatica often improves with time and conservative care, but can persist if nerve compression continues. Recovery from serious injuries such as spinal fractures depends on the level and severity of the injury. A small proportion of people develop chronic back pain, which is best managed with an active, multidisciplinary approach rather than avoidance of activity.

\section*{Prevention and Self-Care}
\begin{itemize}
    \item Exercise regularly, including core strength and back extensor conditioning
    \item Use safe lifting technique: bend the knees, keep the load close, and avoid twisting
    \item Maintain a healthy weight and avoid smoking
    \item Take regular breaks from prolonged sitting; use a supportive chair and workstation
    \item Warm up before sport and progress activity levels gradually
    \item Keep bones strong with adequate calcium, vitamin D, and weight-bearing exercise, especially as you age
    \item See a professional early if back pain is persistent so that it does not become chronic
\end{itemize}

\section*{Sources and Further Reading}
The following sources provide reliable, up-to-date information on back injuries and their treatment.
\begin{itemize}
    \item NHS. \textit{Back pain.} https://www.nhs.uk/conditions/
    \item Mayo Clinic. \textit{Back pain.} https://www.mayoclinic.org
    \item MSD Manual. \textit{Back Pain.} https://www.msdmanuals.com
    \item MedlinePlus. \textit{Back Pain.} https://medlineplus.gov
    \item NICE. \textit{Low back pain and sciatica in over 16s.} https://www.nice.org.uk
    \item World Health Organization. \textit{Low back pain.} https://www.who.int
\end{itemize}
